% !TEX program = xelatex
\documentclass[UTF8,a4paper,11pt,fontset=fandol]{ctexart}
\usepackage[margin=20mm,headheight=15pt]{geometry}
\usepackage{amsmath,amssymb,graphicx,booktabs,longtable,array,tabularx,multirow}
\usepackage{caption,float,enumitem,xcolor,hyperref,fancyhdr}
\definecolor{ink}{HTML}{153A52}
\hypersetup{unicode=true,colorlinks=true,linkcolor=ink,citecolor=ink,urlcolor=ink,pdftitle={血液细胞调控候选的发现与扰动响应预测},pdfauthor={课程分析报告}}
\graphicspath{{figures/}{./}}
\captionsetup{font=small,labelfont=bf,labelsep=quad,skip=6pt}
\linespread{1.10}
\setlength{\parindent}{2em}\setlength{\parskip}{3pt}
\setcounter{tocdepth}{1}
\ctexset{section={format=\Large\bfseries\color{ink}},subsection={format=\normalsize\bfseries\color{ink}}}
\pagestyle{fancy}\fancyhf{}\fancyhead[L]{\small 血液细胞调控候选与扰动响应}\fancyhead[R]{\small 基因组学数据统计分析 · Homework 2}\fancyfoot[C]{\thepage}
\newcommand{\plate}[4]{\begin{figure}[H]\centering\includegraphics[width=\linewidth,height=#4\textheight,keepaspectratio]{#1.pdf}\caption{#3}\label{#2}\end{figure}}
\newcommand{\appage}[1]{\clearpage\subsection{#1}}
\newcommand{\na}{---}
\begin{document}
\thispagestyle{empty}
\begin{center}
{\LARGE\bfseries\color{ink} 血液细胞调控候选的发现\\[5pt]与扰动响应预测}\par\vspace{9pt}
{\large 基于 L1000 表达模式匹配、稳定性筛选与多模型比较}\par\vspace{9pt}
姓名：\underline{\hspace{2.8cm}}\quad 学号：\underline{\hspace{3.5cm}}\par
\small 基因组学数据统计分析 · 第二次作业\quad 2026 年 11 月
\end{center}
\begin{abstract}
从细胞表达状态中寻找值得验证的调控因子，需要把统计匹配与生物学背景联系起来。本研究以造血干细胞为参照，利用 L1000 过表达扰动谱重建血液细胞的表达差异，并以重复抽样确定转录因子候选的优先级。对 21 类目标的筛查显示，扰动字典能够捕获多种血液细胞的差异结构，其中红细胞重建最好，留出基因上的 $R^2$ 为 0.609。深入分析 B 细胞、浆细胞、单核细胞和红细胞后，MYB、负向 BCL6 以及 KLF6／PRDM1 分别形成红系分化、浆细胞程序和髓系状态的机制线索。四类细胞的前十位候选在重新调参的抽样中仍保持至少 90\% 的选择频率，方向全部一致。

进一步在 K562 的单基因及组合 CRISPRa 数据上比较线性基线、GEARS、scGPT 和 scFoundation＋GEARS，将候选发现延伸到干预后的响应预测。共同基因评价中，Ridge 加和与训练条件均值具有接近的平均误差；高变基因评价中，GEARS 的平均误差最低。PRDM1 案例进一步呈现响应幅度与模式相关性两种评价重点。研究建立了以表达程序为起点、以候选稳定性和细胞背景为依据的分析路径，为后续基因调控验证与细胞状态预测提供可追溯的研究线索。
\end{abstract}
\noindent\textbf{关键词：}血液细胞；Lasso；转录因子；扰动预测；单细胞基础模型
\vspace{5pt}
\tableofcontents

\clearpage
\section{引言}
血液系统中不同细胞类型承担免疫、防御与氧运输等功能，其分化伴随着成组基因表达的改变。理解这些差异，需要从大量变化的基因中找到能够组织表达程序的调控因子。平均表达谱呈现了细胞状态，而基因扰动实验提供了改变调控因子后形成的响应模式。把两类信息联系起来，可以将宽泛的细胞差异问题转化为具体的候选发现问题。

L1000 提供了大规模扰动与表达响应的对应关系\cite{R01}。本研究使用课程给定的 MCF7 过表达谱作为模式字典，以目标血液细胞相对造血干细胞（HSC）的差值作为待解释信号。Lasso 从相关的扰动模式中选择联合贡献；TF 注释、重采样和文献证据进一步帮助识别值得研究的因子\cite{R02}。分析的核心问题是：\textbf{哪些调控相关模式能够重建细胞差异，哪些候选在数据变化后仍然保持，哪些生物学机制最值得继续检验？}

研究由候选发现与响应预测两部分组成（图\ref{fig:1}A）。前者覆盖所有非 HSC 类型，重点讨论 B 细胞、浆细胞、单核细胞和红细胞；后者利用独立的 K562 扰动数据，比较多种模型如何预测留出条件下的表达响应。两部分通过共同候选 PRDM1 衔接，从而使基础模型拓展服务于同一研究问题。
\section{结果}
\subsection{表达结构为候选解释提供谱系背景}
主要血液类型在表达空间中形成可辨认的结构，前两个主成分解释约一半的跨类型变异。红细胞与 HSC 的整体差异幅度最大，并呈现血红蛋白和红系相关基因的突出表达；B 细胞与浆细胞则分别具有 B 细胞受体与分泌细胞相关特征（图\ref{fig:1}）。这种结构使四个重点目标具有互补性：既能研究相近谱系之间的状态转变，也能比较功能差异较大的细胞程序。完整高变基因、表达分布及聚类对照见附录\ref{app:A}。
\clearpage
\plate{Fig1}{fig:1}{血液表达结构与研究框架。A：状态、候选与响应预测的联系。B：21 类主要平均谱的 HVG1000 PCA，灰点为其他类型。C：四目标相对 HSC 的差值均方根，覆盖全部 19,994 个去重基因。D：预设谱系标记中的代表基因；热图为逐基因标准化表达，列按 HVG1000 的欧氏距离/Ward 树排序，树与热图共享列顺序。选用标记用于解释谱系特征，完整 29 基因面板和 HVG100 热图在附录 A。}{.80}

\clearpage
\subsection{扰动字典可重建多类血液差异，红细胞最突出}
全部 21 个目标的 Lasso 都优于训练折均值基线。20 个主要目标的留出基因 $R^2$ 中位数为 0.450；四个重点目标中，红细胞达到 0.609，B／浆细胞约为 0.43，单核细胞约为 0.48。图\ref{fig:2}将重建能力、模型规模和替代模型放在同一比较框架中。Ridge 与 Elastic Net 的预测接近 Lasso，支持以 Lasso 的候选解释作为后续主线。采用 1-SE 规则后，模型在保持主要表达结构的同时减少活跃扰动数量。
\plate{Fig2}{fig:2}{表达差异的可重建性。A、B：21 目标的折外 $R^2$ 与全数据模型规模，按最低误差 Lasso 表现排序；mDC 标星，作为附加目标。C：六类模型共享外层划分，颜色区分目标，实心圆为最低误差选择，空心菱形为 1-SE。D：红细胞全部 11,922 个基因的折外预测及边缘分布；该例依据全部主要目标中最高的 Lasso $R^2$ 选取。六边形按基因密度着色，虚线为一致线。完整四目标预测、各指标及调参诊断见附录 B。}{.65}

\clearpage
\subsection{稳定候选呈现细胞特征与共享表达程序}
重复抽样分别得到 B、浆、单核和红细胞的 41、43、52、42 个稳定优先 TF 条目。红系 MYB、浆细胞负向 BCL6 和髓系 KLF6／PRDM1 具有可联系的机制证据；ZNF22 则提供 B 细胞方向的探索线索。各目标候选重叠较少，而 NFYB 和 OVOL2 在四目标中共同出现，显示字典同时捕捉了目标相关与共享的表达模式（图\ref{fig:3}）。
\plate{Fig3}{fig:3}{候选的特异性与共享性。A：按 TF 身份对齐系数分布、选择频率和证据背景；点为 100 次抽样中位数，线为 2.5\%–97.5\% 分位范围。代表因子依据机制联系及共享模式选取；BCL6 在浆细胞中排第 15，因直接文献证据而纳入。B：四目标前十位候选的并集，颜色为选择频率乘以全数据系数符号；B、Plasma、Mono、Ery 分别代表四目标。C：全部稳定候选按规范身份计算 Jaccard。完整前十统计、稳定身份与证据表见附录 C。}{.62}

\clearpage
\subsection{重复选择支持候选优先级，读出范围决定解释层次}
在额外的重新调参抽样中，40 个展示候选的最低选择频率为 90\%，系数中位数方向全部保持。删除候选自身表达读出后，34 个可匹配条目的方向仍一致，最大相对系数变化为 7.05\%。这些结果支持候选贡献由更广的表达程序承载（图\ref{fig:4}A、B）。

读出组成则带来明显变化：使用全部共同基因训练，在留出的 landmark 基因上仍有可辨认的重建结构；把训练也限定在 landmark 后，四目标的 $R^2$ 接近零。等数量非 landmark 对照表现更高，说明信息覆盖与基因组成共同影响候选发现。原尺度差值与逐基因标准化差值对应不同的表达问题，完整对照在附录 C 中分别展示。
\plate{Fig4}{fig:4}{候选与读出的稳健性。A：每目标前十位候选在 20 次重新调参子抽样中的选择频率；点横向轻微错开以显示重合值。B：去除自身读出后的系数相对变化，纵向错位用于区分同目标条目；6 项缺少匹配自身读出，未纳入点图。C：原 Lasso 前十候选在其他分析中的活跃方向保留比例；表达尺度改变时响应定义也随之改变。D：相同的 965 个 landmark 测试基因上，实心方块为 landmark 训练，空心圆为全基因训练，均为最低误差 Lasso。各比较使用相应已冻结结果。}{.62}

\clearpage
\subsection{外部扰动预测揭示读出与泛化条件的作用}
在 Norman 的 K562 CRISPRa 数据中\cite{R11}，条件划分使模型必须预测未参与训练的扰动。共同输出基因上，Ridge 加和与训练条件均值的平均 MSE 十分接近（0.001970 与 0.001973），Ridge 的变化模式相关性更高；在训练集确定的高变基因中，GEARS 平均 MSE 最低（图\ref{fig:5}）。

性能同时取决于训练中是否见过组成基因。Ridge 在两个组成基因均已出现的 seen2 组合上表现突出，而训练条件均值在 seen0 上具有最低平均误差。基础模型与图模型的比较因此需要同时考察读出集合和泛化层次。直接单扰动加和适用于更小的匹配子集，其完整比较保留在附录 D。
\plate{Fig5}{fig:5}{外部扰动预测的完整方法比较。A：每次划分先对条件等权，再计算相对零变化的 MSE 比值，最后对三划分等权平均。圆点与菱形分别为共同基因与训练 HVG1000。B：相同模型顺序下的表达变化 Pearson 相关；零变化的恒定预测无定义相关值。C、D：两种读出下，相对训练条件均值的 MSE 改善百分比，使用共同模型顺序与以零为中心的色标。Single 为未见单扰动；seen0/1/2 表示组合中在最终训练条件里出现的组成基因数。全指标、划分波动、排除直接靶基因及直接加和匹配结果见附录 D。}{.65}

\clearpage
\subsection{PRDM1 将候选身份连接到可观察的扰动响应}
PRDM1 是单核细胞的稳定正向候选，并出现在 Norman 的扰动集合中。依据这一候选身份，在读取 GPU 测试结果前确定其为衔接案例。种子 1 将 PRDM1 单扰动及 CBFA2T3＋PRDM1 组合留作测试，使这一问题能够落到具体响应上（图\ref{fig:6}）。

PRDM1 单扰动的高变基因响应中，GEARS 取得最高相关性，scFoundation＋GEARS 取得最低 MSE。前者较好捕捉变化的相对方向，后者在该例中更贴近整体幅度。这个案例说明，候选发现之后仍需具体比较干预响应，且模式相关与幅度误差能够提供互补信息。
\plate{Fig6}{fig:6}{PRDM1 的跨任务研究衔接。A：血液单核候选身份与 K562 留出扰动之间的联系，两者具有不同细胞背景。B、C：种子 1 的 PRDM1 单扰动和 CBFA2T3＋PRDM1 组合在全部六种可覆盖方法下的 HVG1000 指标，模型顺序一致。D、E：单扰动中分别具有最高相关性和最低 MSE 的两个模型，展示全部 1,000 个评价基因；方法选择依据已说明的两个指标。其余方法及组合的完整指标、预测图见附录 D。该例用于研究同一候选的响应建模，不承担血液背景机制验证。}{.62}

\clearpage
\section{生物学讨论}
\subsection{红系 MYB：把统计首位联系到发育程序}
MYB 在红细胞模型中位列 TF 首位，呈正向贡献，并在所有固定参数子抽样中被选中。其意义来自统计优先级与红系生物学的相互呼应。Bianchi 等在人 CD34$^+$ 造血细胞中发现，MYB 下调影响红系与巨核系命运选择，涉及 KLF1 和 LMO2\cite{R03}；小鼠研究也显示 Myb 参与红系承诺、母细胞阶段推进及终末分裂\cite{R04}。这些结果使 MYB 成为本次候选中最清晰的红系机制线索。

MYB 的作用随发育阶段而变化。因此，当前正向模式匹配更适合指向一个与红系发育相关的表达程序。后续研究可围绕细胞所处阶段比较 MYB 活性、相关表达模块和功能读出，明确哪些状态最能复现当前的联合表达特征。红细胞模型良好的重建表现，也为先在这一谱系中检验分析流程提供了依据。
\subsection{B／浆细胞：方向信息帮助组织机制假设}
浆细胞模型中的 BCL6 是稳定负向候选。已有研究表明，Bcl-6 可直接结合并抑制 Prdm1，Bcl6 缺失促进小鼠浆细胞形成\cite{R05}。因此，BCL6 相关模式的反向贡献与解除分化抑制的机制方向形成联系。把系数符号保留到图表和证据表中，使候选解释能够超越简单的基因名单。

B 细胞中的首位候选 ZNF22 则具有不同的研究定位。本次定向文献检索尚未找到它在 B 细胞分化中的直接机制证据，因而更适合进入独立数据复现和背景验证。与之相对，具有明确 B 细胞发育证据的 EBF1\cite{R08}虽然存在于扰动字典，却在当前联合模型中为零系数。把探索候选与已知谱系因子共同作为研究对象，有助于理解联合表达匹配对细胞背景和状态的选择性。
\subsection{髓系 KLF6／PRDM1：从谱系身份深入功能状态}
单核细胞中的 KLF6 与 PRDM1 均呈稳定正向贡献。Date 等在巨噬细胞中支持 KLF6 参与炎症极化\cite{R06}，随后研究提出 KLF6、PRDM1 与 BCL6 的调控联系\cite{R07}。这为两个候选共同出现提供了具体的功能状态假设：它们可能联系髓系炎症相关表达程序。

这一解释需要保持细胞背景的层次。文献中的巨噬细胞、课程数据中的单核平均谱和外部实验的 K562 分别回答不同问题。最直接的研究延伸是检验独立髓系数据中的同类表达模块，再把候选干预与功能状态读出联系起来。PRDM1 的外部案例增加了响应预测这一层信息，也表明同一因子在不同背景下值得单独建模。
\subsection{共享因子提示共同状态程序}
NFYB 与 OVOL2 在四个目标中分别保持负向和正向稳定贡献，而高变基因中包含较多核糖体相关基因。由此可以提出，部分共享候选可能反映生长、代谢或共同状态程序。后续通过模块注释以及对增殖、应激和分化阶段的分层，可将这类共同变化与谱系相关变化进一步分解。

\clearpage
\subsection{从候选验证连接基因编辑与递送研究}
候选优先级为后续研究提供了明确起点：红系关注 MYB，B／浆细胞关注 BCL6 与探索候选 ZNF22，髓系关注 KLF6／PRDM1。验证可围绕三个连续问题展开：候选是否在独立血液数据中重现；背景明确的干预是否改变对应表达程序；程序改变是否伴随预期细胞状态或功能输出。含供体信息的单细胞数据尤其适合评价跨个体重复性，并帮助分离细胞组成与细胞内状态。

基因编辑技术使转录调控候选能够进一步连接功能元件。Frangoul 等通过编辑自体造血干／祖细胞中 BCL11A 的红系增强子提高胎儿血红蛋白\cite{R09}，说明调控因子的机制理解可以落到细胞特异的调控元件和明确功能输出。对本研究而言，优先任务是建立候选、作用背景和读出之间的一致证据，再讨论合适的调控层级。

递送决定了干预实际发生在哪些细胞、持续多久以及影响怎样的状态分布。靶向 HSC 的 mRNA 脂质纳米颗粒研究\cite{R10}为造血系统的在体调节提供了相关技术背景。对于具有阶段依赖性的因子，靶向性与表达持续时间应与候选身份共同考虑。这里的联系形成研究问题与评价维度，具体实验方案需要依据独立验证后的候选和模型系统制定。
\subsection{基础模型与细胞世界模型：把状态、干预和响应放在一起}
GEARS 将基因关系引入扰动预测\cite{R12}，scGPT 与 scFoundation 则利用大规模单细胞数据学习表达表示\cite{R13,R14}。本研究把这些模型放在统一的条件划分和输出集合中，使比较围绕同一个问题：模型在未见条件下能否预测表达变化。共同基因与高变基因的不同排序说明，评价读出本身决定了我们更关注全局平均结构还是活跃变化。

简单基线在本次比较中具有较强表现。Ridge 在 seen2 上的优势符合其从已见组成基因学习加和效应的结构；完全未见组成基因时，训练条件均值则提供有效的群体参照。近期扰动预测基准同样强调简单方法与严格划分的重要性\cite{R15,R16}。这一结果使后续基础模型研究的重点更加具体：通过背景表示与先验知识，提高新干预条件下的响应预测，同时保留可解释基线作为参照。

从医学世界模型的视角，细胞层面可以构造状态、干预与响应的连续关系：状态由细胞背景及表达程序描述，干预由目标及作用方式描述，响应由转录变化与功能输出描述。本研究覆盖了其中的候选提出与转录响应两个环节。若进一步加入时间、剂量、供体及组织背景，便能够研究状态转移、持续效应与跨背景迁移。将这条路径与递送可达性结合，还可把模型预测转化为更贴近实际研究条件的假设。

\clearpage
\section{材料与方法}
\subsection{数据、注释与探索分析}
课程血液文件包含 20,001 个基因记录和 22 类平均表达谱；MCF7 L1000 过表达矩阵包含 12,328 个基因记录和 2,717 个扰动。重复基因保留首条，分别得到 19,994 与 12,327 个唯一基因；原符号精确交集为 11,940 个。删除在全部血液类型中均为零的 18 个基因后，模型矩阵含 11,922 行。血液数据保持输入尺度，L1000 保留正负分数，所有扰动列均通过数值检查。Human TFs 与 HGNC 注释得到 291 个 TF 扰动条目、290 个规范身份\cite{R02}；不同构建体标签保留独立列。官方 GSE92742 注释识别了 965 个共同 landmark 基因，名称差异通过 Entrez ID 桥接用于注释。

mDC 的血液零值率为 94.00\%，保留建模并单列附加结果。主要探索分析覆盖其余 21 类平均谱，包括 HSC。按跨类型方差选择 HVG100 作热图、HVG1000 作 PCA 和聚类；标准化按基因进行，使用样本标准差。主聚类采用欧氏距离和 Ward 连接，相关距离与平均连接作为对照；改变高变基因数量、尺度和聚类数评价敏感性。
\subsection{回归任务与嵌套验证}
令 $g$ 为基因，$p$ 为扰动，$t$ 为目标细胞，响应及模型为
\begin{equation}
 y_{gt}=E_{gt}-E_{g,\mathrm{HSC}},\qquad
 \widehat y_{gt}=\beta_{0t}+\sum_{p=1}^{2717}X_{gp}\beta_{pt}.
\end{equation}
Lasso 在标准化训练数据上最小化 $\|y-X\beta\|_2^2/(2n)+\lambda\|\beta\|_1$，并拟合截距。观测单位为基因；正负系数表示扰动模式在联合表达重建中的方向。预测还原至响应单位后计算指标。

固定种子 20261107，所有模型共享外层五折、内层三折。训练折独立拟合均值、尺度、降维及参数选择。Lasso 和非负 Lasso 使用 50 个相对惩罚强度，从各训练折的 $\lambda_{\max}$ 到 $10^{-3}\lambda_{\max}$；Elastic Net 混合比例为 0.1、0.5、0.9，各使用同一路径。Ridge 在 $10^{-5}$ 至 $10^{3}$ 搜索 17 个对数间隔惩罚，PCR 成分为 2、5、10、20、50、100，PLS 为 2、5、10、20。21 目标比较 Lasso 与均值基线，四重点目标完成六类模型比较，另对全部目标进行单扰动相关排序。

最低内层 MSE 参数用于预测评价；1-SE 规则从最低误差加一个标准误范围内选择更简约模型。外层每个基因获得一次留出预测，合并计算 $R^2$、RMSE、Pearson 与 Spearman 相关。模型规模来自全数据独立调参后的拟合。候选以全数据 1-SE 模型解释，保持与外层性能评价的用途区分。
\subsection{候选稳定性与敏感性}
每目标进行 100 次无放回 80\% 基因子抽样，每次 9,537 行，固定全数据 1-SE 相对惩罚，并重新估计尺度及 $\lambda_{\max}$。系数绝对值大于 $10^{-6}$ 视为活跃。稳定优先规则为选择频率至少 80\%、活跃抽样中符号一致性至少 90\%、全数据系数非零。按选择频率、系数幅度中位数、全数据幅度依次排序；每目标展示十项。另在相同前 20 次子抽样中重新做三折调参，评价参数选择带来的变化。

敏感性包括 Elastic Net 活跃方向、规范身份候选重叠、相关扰动组、landmark 训练及匹配测试、三组等量非 landmark 对照、逐基因标准化响应和删除候选自身读出。重采样分位线描述条件抽样分布。详细定义、候选表及所有对照见附录 B、C。

\clearpage
\subsection{外部扰动数据与模型输入}
使用 Norman 2019 GSE133344 的 K562 CRISPRa 数据\cite{R11}，采用 scFoundation 发布的含总计数版本（Figshare 24049200）。原数据为 84,986 个细胞、19,264 个基因；因 GEARS GO 映射缺少 IER5L，统一删除相关 843 个细胞及两个条件。保留 84,143 个细胞、225 个条件（含对照）、100 个扰动目标，其中对照细胞为 7,275 个。保留发布者的 log1p 归一化表达，原始总计数独立用于 scFoundation 输入。

统一采用 GEARS simulation 划分，种子为 1、2、3；同条件细胞全部归于同一集合。seen0、seen1、seen2 按最终训练条件中见过的组成基因数重新核对，未见单扰动另组统计。共同输出为 18,977 个基因，每个种子的 HVG1000 由训练数据确定。全部训练目标细胞保留，三个模型共享 64 个固定训练对照作为训练与预测参考池；条件预测取这些参考细胞上的均值，观测值为测试条件所有细胞的均值。
\subsection{模型训练与基线}
GEARS 使用官方模型及 power4＋方向损失，隐藏维数 64。scGPT 加载通用 blood 编码器权重并训练扰动任务模块，核验 153 个编码器张量；每批训练输入 1,536 个基因且包含直接扰动基因，预测覆盖完整共同基因，使用 bfloat16。scFoundation＋GEARS 使用官方 v1 编码器（highres=0）及隐藏维数 512，编码器冻结；64 个参考细胞的 float32 表示缓存经过实时输出一致性检查。三者均在本任务训练集上训练预测部分。

GEARS／scFoundation 使用 Adam，学习率 $10^{-3}$、权重衰减 $5\times10^{-4}$、逐轮 StepLR 系数 0.5、梯度值裁剪 1；scGPT 使用 Adam，学习率 $10^{-4}$、无权重衰减、StepLR 系数 0.9、梯度范数裁剪 1。训练批量分别为 64、24、64，评价批量为 32、12、16（按 GEARS、scFoundation、scGPT 顺序）。GPU 为 A100 80 GB，正式矩阵运算允许 TF32。

用户将预算上限由 20 轮改为 8 轮，验证集连续五轮无改善早停；全部最终最佳检查点均位于前八轮。调整前已完成的额外轮次保留为运行记录，列于附录 D。所有模型遵守共同的划分、输出和评价规则，超参数未依据测试结果选择。

四种基线为零变化、训练条件平均变化、Ridge 加和及单扰动直接加和。后者要求组合的两种单扰动均有训练数据，三个种子分别覆盖 10、16、14 个条件，因此在相同匹配子集内比较所有方法。Ridge 根据训练条件的扰动组成编码学习加和响应；未见组成基因不能从训练中获得其独立效应。
\subsection{评价、文献与复现}
对每个测试条件计算表达变化的 MSE、RMSE、Pearson、Spearman；先按条件等权，再按种子等权汇总。共同基因与训练 HVG1000 分别评价，并补充排除直接扰动基因的指标。零变化等恒定预测的相关系数记为无定义。原型案例覆盖最大、中位、最小真实响应幅度；PRDM1 依据核心候选身份提前确定，随后选取不同评价指标下有解释力的模型图示。

文献以候选、谱系及功能状态检索，核对原始研究的对象与结论；证据表保留直接红系、浆细胞、邻近髓系及探索线索的背景差别。所有分析使用已归档的冻结结果，附录给出必要的完整汇总，LaTeX 工程另存重绘程序和图源。

\clearpage
\section{局限性}
本研究的核心任务是跨基因表达差异重建。血液平均谱缺少供体层面的重复信息，外层留出基因评价的是这一任务的泛化，后续跨供体与跨数据集验证仍有独立价值。基因之间的相关结构也会影响重采样结果，当前稳定性应理解为给定输入与模型规则下的重复选择性。

扰动字典来自 MCF7，目标谱来自血液细胞，外部响应实验来自 K562。系数描述联合表达模式的贡献；将其转化为机制或干预方向，需要结合对应背景中的证据。L1000 的广泛读出包含推断基因，landmark 敏感性表明训练覆盖及读出组成对结果有明显影响。

外部实验包含一个数据集、三次条件划分和八轮预算，基础模型的结论对应本次任务适配与训练设置。scGPT 与 scFoundation 的已公开预训练材料未提供可逐细胞核对的成员清单，与 Norman 数据的样本重叠状态为未核实。直接加和覆盖特定可匹配组合；各种模型的总体结果应与其适用泛化条件共同阅读。
\section{结论}
将血液细胞差异与扰动表达字典联系起来，可以形成从状态结构、联合重建到稳定候选的完整证据链。红细胞提供了最清晰的重建实例，MYB、BCL6 和 KLF6／PRDM1 则把统计优先级连接到可检验的谱系与功能状态问题。

外部扰动比较进一步表明，响应预测需要同时考虑读出范围和训练中已见的干预组成。PRDM1 案例把候选身份落实为具体响应问题，也展示了幅度误差与模式相关性的互补作用。这套分析为后续独立血液数据验证、背景明确的基因调控研究，以及状态—干预—响应模型提供了可复用起点。

\vspace{1em}
\noindent\textbf{学生补充说明：}AI 智能体配置与使用方式、本人检查结果的具体例子，以及个人研究联系和反思，请在配套工作记录的预留部分填写。报告署名请填写首页姓名与学号。

\clearpage
\phantomsection
\addcontentsline{toc}{section}{参考文献}
\begin{thebibliography}{99}\small
\bibitem{R01} Subramanian A;  Narayan R;  Corsello SM; et al.. A Next Generation Connectivity Map: L1000 Platform and the First 1,000,000 Profiles. Cell, 2017. DOI: \nolinkurl{10.1016/j.cell.2017.10.049}.
\bibitem{R02} Lambert SA;  Jolma A;  Campitelli LF; et al.. The Human Transcription Factors. Cell, 2018. DOI: \nolinkurl{10.1016/j.cell.2018.01.029}.
\bibitem{R03} Bianchi E;  Zini R;  Salati S; et al.. c-myb supports erythropoiesis through the transactivation of KLF1 and LMO2 expression. Blood, 2010. DOI: \nolinkurl{10.1182/blood-2009-08-238311}.
\bibitem{R04} Vegiopoulos A;  García P;  Emambokus N; et al.. Coordination of erythropoiesis by the transcription factor c-Myb. Blood, 2006. DOI: \nolinkurl{10.1182/blood-2005-07-2968}.
\bibitem{R05} Tunyaplin C;  Shaffer AL;  Angelin-Duclos CD; et al.. Direct repression of prdm1 by Bcl-6 inhibits plasmacytic differentiation. Journal of immunology (Baltimore, Md. : 1950), 2004. DOI: \nolinkurl{10.4049/jimmunol.173.2.1158}.
\bibitem{R06} Date D;  Das R;  Narla G; et al.. Kruppel-like transcription factor 6 regulates inflammatory macrophage polarization. The Journal of biological chemistry, 2014. DOI: \nolinkurl{10.1074/jbc.M113.526749}.
\bibitem{R07} Kim GD;  Das R;  Goduni L; et al.. Kruppel-like Factor 6 Promotes Macrophage-mediated Inflammation by Suppressing B Cell Leukemia/Lymphoma 6 Expression. The Journal of biological chemistry, 2016. DOI: \nolinkurl{10.1074/jbc.M116.738617}.
\bibitem{R08} Lin H;  Grosschedl R. Failure of B-cell differentiation in mice lacking the transcription factor EBF. Nature, 1995. DOI: \nolinkurl{10.1038/376263a0}.
\bibitem{R09} Frangoul H;  Altshuler D;  Cappellini MD; et al.. CRISPR-Cas9 Gene Editing for Sickle Cell Disease and $\beta$-Thalassemia. The New England journal of medicine, 2021. DOI: \nolinkurl{10.1056/NEJMoa2031054}.
\bibitem{R10} Breda L;  Papp TE;  Triebwasser MP; et al.. In vivo hematopoietic stem cell modification by mRNA delivery. Science (New York, N.Y.), 2023. DOI: \nolinkurl{10.1126/science.ade6967}.
\bibitem{R11} Norman TM;  Horlbeck MA;  Replogle JM; et al.. Exploring genetic interaction manifolds constructed from rich single-cell phenotypes. Science, 2019. DOI: \nolinkurl{10.1126/science.aax4438}.
\bibitem{R12} Roohani Y;  Huang K;  Leskovec J. Predicting transcriptional outcomes of novel multigene perturbations with GEARS. Nature biotechnology, 2024. DOI: \nolinkurl{10.1038/s41587-023-01905-6}.
\bibitem{R13} Cui H;  Wang C;  Maan H; et al.. scGPT: toward building a foundation model for single-cell multi-omics using generative AI. Nature methods, 2024. DOI: \nolinkurl{10.1038/s41592-024-02201-0}.
\bibitem{R14} Hao M;  Gong J;  Zeng X; et al.. Large-scale foundation model on single-cell transcriptomics. Nature methods, 2024. DOI: \nolinkurl{10.1038/s41592-024-02305-7}.
\bibitem{R15} Ahlmann-Eltze C;  Huber W;  Anders S. Deep-learning-based gene perturbation effect prediction does not yet outperform simple linear baselines. Nature Methods, 2025. DOI: \nolinkurl{10.1038/s41592-025-02772-6}.
\bibitem{R16} Wenkel F;  Tu W;  Masschelein C; et al.. TxPert: using multiple knowledge graphs for prediction of transcriptomic perturbation effects. Nature Biotechnology, 2026. DOI: \nolinkurl{10.1038/s41587-026-03113-4}.
\end{thebibliography}
\clearpage
\appendix
\clearpage\section{数据与探索分析}\label{app:A}\setcounter{figure}{0}\setcounter{table}{0}\renewcommand{\thefigure}{A\arabic{figure}}\renewcommand{\thetable}{A\arabic{table}}
本附录给出原始矩阵质控、细胞覆盖及完整探索图。除 mDC 的敏感性分析外，主要 EDA 覆盖 HSC 和 20 个目标。
{\small\linespread{1}\selectfont\setlength{\tabcolsep}{4pt}\renewcommand{\arraystretch}{1.08}\begin{longtable}{lrrrrr}\caption{原始输入矩阵的质量与尺度。缺失和无穷值均为零。}\\\toprule 输入 & 行数 & 列数 & 删除重复行 & 最小值 & 最大值\\\midrule\endfirsthead\toprule 输入 & 行数 & 列数 & 删除重复行 & 最小值 & 最大值\\\midrule\endhead
blood\_raw & 20001 & 22 & 7 & 0.000 & 8.778\\
L1000\_raw & 12328 & 2717 & 1 & -10.000 & 10.000\\
\bottomrule\end{longtable}}
血液去重后 19,994 个基因，L1000 去重后 12,327 个。模型交集去除全血液零值行后为 11,922 个基因；全部 2,717 个扰动列保留，291 个为 TF 条目。重复处理逐条保留审计，矩阵按原符号对齐，HGNC 映射用于注释。
\plate{02_F1_expression_distributions}{a:dist}{表达覆盖与分布。A：全部 22 类型的零值比例。B：HSC 与重点目标的跨基因累积分布；symlog 为显示坐标，输入表达尺度保持。}{0.56}
\clearpage\subsection{完整高变基因表达图}
\plate{Supp_HVG100}{a:hvg}{完整 HVG100 逐基因 z-score。A、B 按原 Ward 行聚类顺序分别显示前、后 50 个基因；两块共享同一细胞列顺序和色标。列顺序来自图示 HVG100 的 Ward 聚类；完整列树见顶部。图中 56 个基因带有核糖体相关分组注释。}{0.83}
\clearpage\subsection{主成分与谱系差异的补充视图}
\plate{02_F3_PCA}{a:pca}{完整 PCA：A 为 PC1/PC2，B 为 PC1/PC3，C 为解释方差。前两轴累计 51.54\%，前三轴 63.67\%。所有点为细胞类型平均谱。}{0.33}
\plate{02_F4_HSC_differences}{a:markers}{相对 HSC 的差异。A 为四目标差值 RMS；B 为预设 29 个谱系标记在逐基因标准化后与 HSC 的差值，保留全部可用标记。}{0.44}
\clearpage\subsection{聚类距离与 mDC 的影响}
\plate{02_F5_clustering_comparison}{a:cluster}{HVG1000 下欧氏距离/Ward 与 Pearson 相关距离/平均连接的对照。HSC 在两者中均最接近造血前体。}{0.37}
\plate{02_S1_mDC_sensitivity}{a:mdc}{mDC 敏感性：A 将 mDC 投影到固定主要 PCA 空间；B 将全部 22 类型重新拟合。mDC 的高零值率及其对全局坐标的影响单独显示。}{0.33}
\clearpage\subsection{全部细胞的覆盖与聚类稳定性}
{\small\linespread{1}\selectfont\setlength{\tabcolsep}{4pt}\renewcommand{\arraystretch}{1.08}\begin{longtable}{lrrr}\caption{细胞覆盖。零值率以去重后的 19,994 个基因为分母。}\\\toprule 细胞类型 & 零值率（\%） & 均值 & 最大值\\\midrule\endfirsthead\toprule 细胞类型 & 零值率（\%） & 均值 & 最大值\\\midrule\endhead
HSC & 29.2 & 0.173 & 4.565\\
CMP & 55.9 & 0.190 & 4.436\\
Erythrocyte & 20.2 & 0.019 & 8.778\\
Platelet & 39.5 & 0.093 & 5.491\\
Macrophage & 19.2 & 0.163 & 4.374\\
B cell & 10.8 & 0.145 & 4.189\\
Monocyte & 2.7 & 0.164 & 5.044\\
NK cell & 14.8 & 0.164 & 4.368\\
CD4 T & 5.4 & 0.142 & 4.117\\
CD8 T & 11.6 & 0.164 & 4.374\\
Basophil & 45.4 & 0.149 & 3.981\\
Neutrophil & 3.6 & 0.119 & 5.215\\
mDC & 94.0 & 0.060 & 6.317\\
pDC & 37.4 & 0.194 & 4.196\\
Plasma cell & 16.0 & 0.149 & 4.413\\
NK T & 29.1 & 0.165 & 4.274\\
Treg & 48.0 & 0.156 & 4.635\\
Classical mono. & 3.9 & 0.177 & 4.964\\
Non-classical mono. & 26.0 & 0.186 & 4.778\\
Naive CD4 T & 24.2 & 0.145 & 4.526\\
Intermediate mono. & 19.4 & 0.195 & 4.517\\
Precursor & 27.2 & 0.147 & 4.309\\
\bottomrule\end{longtable}}
\plate{02_S2_clustering_sensitivity}{a:clustsens}{聚类敏感性。改变 HVG 数目、表达尺度和距离定义，在聚类数 3–8 下以 ARI 对照主聚类；ARI 衡量分组一致性。}{0.28}
\clearpage\section{统计模型与拟合诊断}\label{app:B}\setcounter{figure}{0}\setcounter{table}{0}\renewcommand{\thefigure}{B\arabic{figure}}\renewcommand{\thetable}{B\arabic{table}}
本附录的 r 和 $\rho$ 分别为 Pearson 与 Spearman 相关。最低误差模型用于预测比较，1-SE 用于简约候选解释；模型规模均来自全数据最终拟合。
{\small\linespread{1}\selectfont\setlength{\tabcolsep}{4pt}\renewcommand{\arraystretch}{1.08}\begin{longtable}{lrrrrrr}\caption{全部目标 Lasso 的折外性能与模型规模。min/1SE 列保留两套选择。}\\\toprule 目标 & $R^2_{min}$ & $R^2_{1SE}$ & RMSE$_{min}$ & $r_{min}$ & $\rho_{min}$ & 活跃数 min/1SE\\\midrule\endfirsthead\toprule 目标 & $R^2_{min}$ & $R^2_{1SE}$ & RMSE$_{min}$ & $r_{min}$ & $\rho_{min}$ & 活跃数 min/1SE\\\midrule\endhead
Erythrocyte & 0.609 & 0.568 & 0.2413 & 0.781 & 0.695 & 1142/593\\
Neutrophil & 0.547 & 0.533 & 0.2954 & 0.740 & 0.603 & 991/708\\
Macrophage & 0.538 & 0.516 & 0.1979 & 0.734 & 0.631 & 1019/824\\
Intermediate mono. & 0.504 & 0.475 & 0.2078 & 0.710 & 0.635 & 1038/787\\
Classical mono. & 0.494 & 0.460 & 0.2019 & 0.704 & 0.634 & 1041/733\\
CD8 T & 0.486 & 0.463 & 0.1630 & 0.698 & 0.606 & 1117/890\\
Monocyte & 0.483 & 0.445 & 0.1927 & 0.696 & 0.628 & 1054/683\\
Non-classical mono. & 0.472 & 0.451 & 0.1946 & 0.688 & 0.605 & 1108/768\\
NK cell & 0.470 & 0.447 & 0.1757 & 0.686 & 0.580 & 1102/842\\
CD4 T & 0.464 & 0.435 & 0.1402 & 0.681 & 0.627 & 1099/677\\
Platelet & 0.437 & 0.406 & 0.3251 & 0.662 & 0.546 & 974/798\\
Basophil & 0.430 & 0.406 & 0.2416 & 0.656 & 0.512 & 838/652\\
Plasma cell & 0.429 & 0.393 & 0.1620 & 0.655 & 0.583 & 1111/784\\
B cell & 0.427 & 0.397 & 0.1595 & 0.654 & 0.611 & 1089/656\\
NK T & 0.409 & 0.391 & 0.1657 & 0.640 & 0.557 & 976/681\\
Naive CD4 T & 0.401 & 0.369 & 0.1539 & 0.634 & 0.592 & 1051/612\\
Treg & 0.400 & 0.371 & 0.1879 & 0.633 & 0.540 & 1011/558\\
pDC & 0.393 & 0.365 & 0.1475 & 0.628 & 0.547 & 1043/803\\
Precursor & 0.312 & 0.284 & 0.0898 & 0.560 & 0.529 & 953/653\\
mDC & 0.171 & 0.108 & 0.4654 & 0.414 & 0.478 & 435/226\\
CMP & 0.127 & 0.085 & 0.1601 & 0.357 & 0.286 & 580/204\\
\bottomrule\end{longtable}}
{\footnotesize\linespread{1}\selectfont\setlength{\tabcolsep}{4pt}\renewcommand{\arraystretch}{1.08}\begin{longtable}{lrrrrr}\caption{1-SE 与均值基线的补充指标。基线由各训练折均值构成。}\\\toprule 目标 & RMSE$_{1SE}$ & $r_{1SE}$ & $\rho_{1SE}$ & $R^2_{mean}$ & RMSE$_{mean}$\\\midrule\endfirsthead\toprule 目标 & RMSE$_{1SE}$ & $r_{1SE}$ & $\rho_{1SE}$ & $R^2_{mean}$ & RMSE$_{mean}$\\\midrule\endhead
Erythrocyte & 0.2536 & 0.754 & 0.675 & -0.0001 & 0.3861\\
Neutrophil & 0.2997 & 0.730 & 0.588 & -0.0001 & 0.4388\\
Macrophage & 0.2025 & 0.718 & 0.611 & -0.0001 & 0.2910\\
Intermediate mono. & 0.2137 & 0.690 & 0.614 & -0.0002 & 0.2951\\
Classical mono. & 0.2086 & 0.679 & 0.611 & -0.0002 & 0.2839\\
CD8 T & 0.1667 & 0.681 & 0.590 & -0.0001 & 0.2274\\
Monocyte & 0.1998 & 0.668 & 0.605 & -0.0002 & 0.2681\\
Non-classical mono. & 0.1986 & 0.672 & 0.592 & -0.0002 & 0.2680\\
NK cell & 0.1795 & 0.669 & 0.564 & -0.0001 & 0.2414\\
CD4 T & 0.1439 & 0.660 & 0.607 & -0.0001 & 0.1914\\
Platelet & 0.3340 & 0.637 & 0.534 & -0.0002 & 0.4333\\
Basophil & 0.2467 & 0.637 & 0.487 & -0.0002 & 0.3200\\
Plasma cell & 0.1670 & 0.628 & 0.557 & -0.0002 & 0.2143\\
B cell & 0.1636 & 0.631 & 0.594 & -0.0000 & 0.2107\\
NK T & 0.1682 & 0.626 & 0.546 & -0.0002 & 0.2155\\
Naive CD4 T & 0.1580 & 0.609 & 0.570 & -0.0002 & 0.1988\\
Treg & 0.1924 & 0.610 & 0.525 & -0.0001 & 0.2426\\
pDC & 0.1509 & 0.604 & 0.523 & -0.0002 & 0.1894\\
Precursor & 0.0917 & 0.533 & 0.516 & -0.0002 & 0.1083\\
mDC & 0.4829 & 0.339 & 0.461 & -0.0002 & 0.5113\\
CMP & 0.1639 & 0.300 & 0.240 & -0.0002 & 0.1714\\
\bottomrule\end{longtable}}
\clearpage\subsection{六类统计模型的完整比较}
{\footnotesize\linespread{1}\selectfont\setlength{\tabcolsep}{4pt}\renewcommand{\arraystretch}{1.08}\begin{longtable}{lrrrrrr}\caption{四目标所有模型与选择规则。各行均采用相同外层基因划分。}\\\toprule 目标 & 模型 & 选择 & $R^2$ & RMSE & $r$ & $\rho$\\\midrule\endfirsthead\toprule 目标 & 模型 & 选择 & $R^2$ & RMSE & $r$ & $\rho$\\\midrule\endhead
B & Lasso & min & 0.427 & 0.1595 & 0.654 & 0.611\\
B & Lasso & 1se & 0.397 & 0.1636 & 0.631 & 0.594\\
B & Elastic Net & min & 0.431 & 0.1590 & 0.657 & 0.614\\
B & Elastic Net & 1se & 0.404 & 0.1627 & 0.636 & 0.598\\
B & Ridge & min & 0.435 & 0.1584 & 0.660 & 0.617\\
B & Ridge & 1se & 0.418 & 0.1607 & 0.647 & 0.611\\
B & Nonneg. Lasso & min & 0.169 & 0.1922 & 0.411 & 0.456\\
B & Nonneg. Lasso & 1se & 0.133 & 0.1963 & 0.375 & 0.430\\
B & PCR & min & 0.151 & 0.1941 & 0.389 & 0.398\\
B & PCR & 1se & 0.138 & 0.1956 & 0.372 & 0.387\\
B & PLS & min & 0.414 & 0.1613 & 0.652 & 0.608\\
B & PLS & 1se & 0.393 & 0.1641 & 0.635 & 0.595\\
B & Mean & baseline & -0.000 & 0.2107 & -0.008 & -0.004\\
Plasma & Lasso & min & 0.429 & 0.1620 & 0.655 & 0.583\\
Plasma & Lasso & 1se & 0.393 & 0.1670 & 0.628 & 0.557\\
Plasma & Elastic Net & min & 0.436 & 0.1610 & 0.661 & 0.587\\
Plasma & Elastic Net & 1se & 0.401 & 0.1659 & 0.634 & 0.562\\
Plasma & Ridge & min & 0.443 & 0.1600 & 0.666 & 0.592\\
Plasma & Ridge & 1se & 0.422 & 0.1630 & 0.651 & 0.580\\
Plasma & Nonneg. Lasso & min & 0.127 & 0.2002 & 0.357 & 0.357\\
Plasma & Nonneg. Lasso & 1se & 0.078 & 0.2057 & 0.305 & 0.284\\
Plasma & PCR & min & 0.136 & 0.1992 & 0.369 & 0.326\\
Plasma & PCR & 1se & 0.110 & 0.2022 & 0.331 & 0.289\\
Plasma & PLS & min & 0.419 & 0.1633 & 0.656 & 0.582\\
Plasma & PLS & 1se & 0.403 & 0.1656 & 0.641 & 0.562\\
Plasma & Mean & baseline & -0.000 & 0.2143 & -0.017 & -0.018\\
Mono & Lasso & min & 0.483 & 0.1927 & 0.696 & 0.628\\
Mono & Lasso & 1se & 0.445 & 0.1998 & 0.668 & 0.605\\
Mono & Elastic Net & min & 0.487 & 0.1919 & 0.699 & 0.632\\
Mono & Elastic Net & 1se & 0.450 & 0.1989 & 0.672 & 0.608\\
Mono & Ridge & min & 0.496 & 0.1904 & 0.704 & 0.638\\
Mono & Ridge & 1se & 0.478 & 0.1936 & 0.692 & 0.631\\
Mono & Nonneg. Lasso & min & 0.198 & 0.2401 & 0.445 & 0.449\\
Mono & Nonneg. Lasso & 1se & 0.152 & 0.2469 & 0.407 & 0.415\\
Mono & PCR & min & 0.232 & 0.2350 & 0.482 & 0.421\\
Mono & PCR & 1se & 0.216 & 0.2374 & 0.465 & 0.406\\
Mono & PLS & min & 0.476 & 0.1941 & 0.696 & 0.624\\
Mono & PLS & 1se & 0.441 & 0.2004 & 0.668 & 0.596\\
Mono & Mean & baseline & -0.000 & 0.2681 & -0.021 & -0.008\\
Ery & Lasso & min & 0.609 & 0.2413 & 0.781 & 0.695\\
Ery & Lasso & 1se & 0.568 & 0.2536 & 0.754 & 0.675\\
Ery & Elastic Net & min & 0.613 & 0.2400 & 0.783 & 0.698\\
Ery & Elastic Net & 1se & 0.572 & 0.2524 & 0.757 & 0.677\\
Ery & Ridge & min & 0.615 & 0.2396 & 0.784 & 0.700\\
Ery & Ridge & 1se & 0.587 & 0.2481 & 0.767 & 0.692\\
Ery & Nonneg. Lasso & min & 0.336 & 0.3146 & 0.579 & 0.568\\
Ery & Nonneg. Lasso & 1se & 0.298 & 0.3235 & 0.554 & 0.553\\
Ery & PCR & min & 0.390 & 0.3016 & 0.624 & 0.537\\
Ery & PCR & 1se & 0.372 & 0.3059 & 0.610 & 0.524\\
Ery & PLS & min & 0.599 & 0.2444 & 0.776 & 0.687\\
Ery & PLS & 1se & 0.573 & 0.2522 & 0.759 & 0.668\\
Ery & Mean & baseline & -0.000 & 0.3861 & -0.014 & -0.000\\
\bottomrule\end{longtable}}
\clearpage\subsection{全部重点目标的留出预测}
\plate{03_F2_focus_OOF_predictions}{b:oof}{四目标最低误差 Lasso 折外预测。每个基因仅使用其所属测试折的预测，全部 11,922 个基因均进入图中。六边形颜色为基因数的对数色阶，各面板独立；虚线为一致线。}{0.78}
预测在低幅度区域集中，大幅变化呈现收缩。评价以整套折外预测合并计算，因而图示代表所有共同基因的重建情况。模型预测单位与目标相对 HSC 的输入尺度差值相同。
\clearpage\subsection{调参曲线与单次拟合初筛}
\plate{03_F3_tuning_curves}{b:tuning}{全数据三折调参曲线：均值 MSE 及一个标准误，蓝虚线与朱红点线分别为最低误差与 1-SE。该图解释最终候选模型选择，独立外层评价见表 B1。}{0.43}
{\footnotesize\linespread{1}\selectfont\setlength{\tabcolsep}{4pt}\renewcommand{\arraystretch}{1.08}\begin{longtable}{rllll}\caption{单次全数据 1-SE 拟合按系数幅度的前十 TF。括号为标准化系数，排序与稳定性排序用途不同。}\\\toprule 名次 & B & Plasma & Mono & Ery\\\midrule\endfirsthead\toprule 名次 & B & Plasma & Mono & Ery\\\midrule\endhead
1 & ZNF22 (+0.069) & OVOL2 (+0.072) & PAX8 (+0.088) & MYB (+0.078)\\
2 & RORC (+0.045) & (DDIT3) (+0.067) & KLF6 (+0.084) & NFYB (-0.053)\\
3 & ZNF296 (+0.040) & ZNF653 (+0.067) & TCF7L2 (+0.067) & IRF3 (+0.050)\\
4 & ZNF174 (-0.040) & NCOA3 (-0.064) & TFAP2A (+0.067) & PREB (+0.047)\\
5 & STAT4 (+0.037) & LHX8 (-0.059) & OVOL2 (+0.062) & ELK3 (+0.041)\\
6 & KLF3 (+0.036) & ERF (+0.053) & HIF1A (+0.060) & CREB3 (+0.040)\\
7 & NFYB (-0.035) & CLOCK (-0.049) & IRF3 (+0.055) & NCOA3 (-0.040)\\
8 & ZNF114 (+0.035) & TFAP2A (+0.047) & SIM2 (-0.050) & RARA (-0.037)\\
9 & HOXB5 (-0.034) & BCL6 (-0.041) & PRDM1 (+0.048) & KLF6 (+0.036)\\
10 & DLX3 (+0.034) & TP53 (-0.040) & ZNF546 (+0.044) & ZNF790 (-0.033)\\
\bottomrule\end{longtable}}
\clearpage\subsection{完整系数路径的代表变量}
\plate{03_F4_coefficient_paths}{b:paths}{每目标全数据 1-SE 模型中幅度最大的十二个活跃扰动，显示完整正则化路径。横轴由强到弱惩罚，蓝线为最低误差，朱红线为 1-SE。正负号贯穿路径保留，变量标签为原始扰动标签。}{0.78}
全部 2,717 个扰动参与拟合。单扰动相关性按与目标差值的绝对 Pearson 相关排序，用于直观参照；联合回归系数还取决于其他扰动模式，故两种排序分别保留。求解过程记录了收敛及边界参数，冻结结果通过数值验收。
\clearpage\section{候选与稳健性证据}\label{app:C}\setcounter{figure}{0}\setcounter{table}{0}\renewcommand{\thefigure}{C\arabic{figure}}\renewcommand{\thetable}{C\arabic{table}}
B、Plasma、Mono、Ery 对应四目标。频率为选中比例；固定抽样为 100 次，重新调参为 20 次。所有展示条目的固定选择频率及活跃符号一致性均为 100\%。分位数描述子抽样系数分布。
{\footnotesize\linespread{1}\selectfont\setlength{\tabcolsep}{4pt}\renewcommand{\arraystretch}{1.08}\begin{longtable}{lrlrrlrr}\caption{40 个展示候选的完整统计及自身读出删除结果。最后一列为相对系数变化百分比，--- 为没有匹配自身读出。}\\\toprule 目标 & 排名 & TF & 全数据 & 中位数 & 2.5\%–97.5\% & 重选\% & 删除变化\%\\\midrule\endfirsthead\toprule 目标 & 排名 & TF & 全数据 & 中位数 & 2.5\%–97.5\% & 重选\% & 删除变化\%\\\midrule\endhead
B & 1 & ZNF22 & 0.069 & 0.068 & [0.041, 0.090] & 100 & 0.08\\
B & 2 & DLX3 & 0.034 & 0.039 & [0.018, 0.067] & 100 & ---\\
B & 3 & RORC & 0.045 & 0.039 & [0.020, 0.059] & 100 & 0.21\\
B & 4 & ZNF174 & -0.040 & -0.037 & [-0.051, -0.019] & 100 & -0.01\\
B & 5 & STAT4 & 0.037 & 0.037 & [0.015, 0.058] & 100 & -0.45\\
B & 6 & ZNF296 & 0.040 & 0.036 & [0.010, 0.059] & 100 & ---\\
B & 7 & KLF3 & 0.036 & 0.034 & [0.011, 0.067] & 90 & 0.03\\
B & 8 & HOXB5 & -0.034 & -0.033 & [-0.048, -0.015] & 100 & -0.08\\
B & 9 & NFYB & -0.035 & -0.032 & [-0.053, -0.010] & 100 & 0.03\\
B & 10 & ZNF114 & 0.035 & 0.032 & [0.018, 0.052] & 100 & ---\\
Plasma & 1 & OVOL2 & 0.072 & 0.065 & [0.047, 0.087] & 100 & -0.01\\
Plasma & 2 & NCOA3 & -0.064 & -0.065 & [-0.087, -0.046] & 100 & -7.05\\
Plasma & 3 & ZNF653 & 0.067 & 0.063 & [0.040, 0.083] & 100 & ---\\
Plasma & 4 & (DDIT3) & 0.067 & 0.057 & [0.037, 0.080] & 100 & -0.06\\
Plasma & 5 & LHX8 & -0.059 & -0.047 & [-0.069, -0.026] & 100 & ---\\
Plasma & 6 & ERF & 0.053 & 0.046 & [0.022, 0.072] & 100 & 0.03\\
Plasma & 7 & CLOCK & -0.049 & -0.041 & [-0.061, -0.029] & 100 & 0.12\\
Plasma & 8 & ZNF558 & 0.039 & 0.039 & [0.018, 0.058] & 100 & ---\\
Plasma & 9 & TFAP2A & 0.047 & 0.038 & [0.017, 0.064] & 100 & 0.09\\
Plasma & 10 & TP53 & -0.040 & -0.037 & [-0.064, -0.014] & 100 & 0.91\\
Mono & 1 & PAX8 & 0.088 & 0.087 & [0.053, 0.110] & 95 & 0.24\\
Mono & 2 & KLF6 & 0.084 & 0.082 & [0.064, 0.098] & 100 & -1.21\\
Mono & 3 & TCF7L2 & 0.067 & 0.065 & [0.050, 0.079] & 100 & 0.06\\
Mono & 4 & TFAP2A & 0.067 & 0.064 & [0.042, 0.087] & 100 & -0.01\\
Mono & 5 & OVOL2 & 0.062 & 0.060 & [0.039, 0.077] & 100 & 0.00\\
Mono & 6 & IRF3 & 0.055 & 0.057 & [0.037, 0.080] & 100 & 0.09\\
Mono & 7 & HIF1A & 0.060 & 0.055 & [0.040, 0.072] & 100 & -0.17\\
Mono & 8 & PRDM1 & 0.048 & 0.051 & [0.031, 0.070] & 100 & -0.05\\
Mono & 9 & SIM2 & -0.050 & -0.049 & [-0.070, -0.032] & 100 & 0.37\\
Mono & 10 & SATB2 & 0.040 & 0.043 & [0.006, 0.065] & 95 & -0.01\\
Ery & 1 & MYB & 0.078 & 0.071 & [0.038, 0.097] & 100 & 1.30\\
Ery & 2 & IRF3 & 0.050 & 0.052 & [0.011, 0.073] & 100 & 0.09\\
Ery & 3 & NFYB & -0.053 & -0.052 & [-0.070, -0.026] & 100 & 0.00\\
Ery & 4 & PREB & 0.047 & 0.047 & [0.034, 0.060] & 100 & -0.02\\
Ery & 5 & ELK3 & 0.041 & 0.044 & [0.019, 0.058] & 100 & 0.05\\
Ery & 6 & NCOA3 & -0.040 & -0.039 & [-0.058, -0.025] & 100 & 0.34\\
Ery & 7 & CREB3 & 0.040 & 0.038 & [0.014, 0.051] & 100 & -0.12\\
Ery & 8 & RARA & -0.037 & -0.037 & [-0.053, -0.010] & 100 & 0.50\\
Ery & 9 & OVOL2 & 0.032 & 0.032 & [0.014, 0.046] & 100 & 0.07\\
Ery & 10 & KLF6 & 0.036 & 0.032 & [0.019, 0.044] & 100 & 1.79\\
\bottomrule\end{longtable}}
\clearpage\subsection{全部稳定身份与交集}
\noindent\textbf{B cell（41 个规范身份）：} BMAL1、DDIT3、DLX3、EGR1、ESR1、FOXR1、HOXB5、IRF5、JUNB、KLF3、MAFG、NFIL3、NFYB、OSR2、OTX1、OTX2、OVOL2、RFX4、RFX6、RORC、SALL4、SNAI3、SREBF1、STAT4、TFAP4、ZBTB20、ZNF114、ZNF136、ZNF174、ZNF22、ZNF296、ZNF317、ZNF32、ZNF331、ZNF426、ZNF449、ZNF473、ZNF621、ZNF653、ZNF747、ZNF785。\par\medskip
\noindent\textbf{Plasma cell（43 个规范身份）：} BCL6、CEBPG、CLOCK、DDIT3、EBF1、ERF、GATA2、HIF1A、IRF3、IRF5、ISX、KAT7、LHX8、MAX、MBNL2、MSX2、MTERF3、MTF2、NCOA3、NFIB、NFKB1、NFYB、NR1H3、OVOL2、SATB2、SOX2、STAT1、STAT3、STAT5A、TCF12、TFAP2A、TP53、WT1、ZNF134、ZNF136、ZNF274、ZNF398、ZNF548、ZNF555、ZNF558、ZNF653、ZNF747、ZNF791。\par\medskip
\noindent\textbf{Monocyte（52 个规范身份）：} CDX2、CEBPG、CREB1、CTCF、ELF1、ERF、FOXP3、HIF1A、HNF4A、HOXB13、HOXC10、IRF2、IRF3、ISX、KLF6、MAFG、MTERF3、NFKB1、NFYB、OSR2、OVOL2、PAX8、PRDM1、SATB1、SATB2、SIM2、SOX2、STAT2、TCF7L2、TFAP2A、TFAP4、TGIF2LY、TP53、VDR、ZBTB20、ZBTB25、ZBTB26、ZBTB49、ZIM3、ZNF112、ZNF114、ZNF174、ZNF300、ZNF490、ZNF546、ZNF555、ZNF558、ZNF576、ZNF583、ZNF597、ZNF677、ZNF785。\par\medskip
\noindent\textbf{Erythrocyte（42 个规范身份）：} CDX2、CREB3、CREB3L1、CREB5、EBF1、EGR1、ELK3、HMGA1、HOXA6、HOXB7、HOXC9、IRF3、KLF6、LHX8、MAX、MEIS2、MTERF3、MYB、NCOA3、NFYB、NR5A2、OVOL2、PRDM1、PREB、RARA、RBPJ、SREBF1、STAT5A、TCF12、TRAFD1、ZBED1、ZBTB49、ZGPAT、ZNF200、ZNF268、ZNF300、ZNF555、ZNF57、ZNF581、ZNF626、ZNF77、ZNF790。\par\medskip
{\footnotesize\linespread{1}\selectfont\setlength{\tabcolsep}{4pt}\renewcommand{\arraystretch}{1.08}\begin{longtable}{lrrp{10cm}}\caption{成对交集。构建体按规范身份合并后计算。}\\\toprule 目标对 & 交集数 & Jaccard & 共享身份\\\midrule\endfirsthead\toprule 目标对 & 交集数 & Jaccard & 共享身份\\\midrule\endhead
B / Plasma & 7 & 0.091 & DDIT3, IRF5, NFYB, OVOL2, ZNF136, ZNF653, ZNF747\\
B / Mono & 9 & 0.107 & MAFG, NFYB, OSR2, OVOL2, TFAP4, ZBTB20, ZNF114, ZNF174, ZNF785\\
B / Ery & 4 & 0.051 & EGR1, NFYB, OVOL2, SREBF1\\
Plasma / Mono & 15 & 0.188 & CEBPG, ERF, HIF1A, IRF3, ISX, MTERF3, NFKB1, NFYB, OVOL2, SATB2, SOX2, TFAP2A, TP53, ZNF555, ZNF558\\
Plasma / Ery & 11 & 0.149 & EBF1, IRF3, LHX8, MAX, MTERF3, NCOA3, NFYB, OVOL2, STAT5A, TCF12, ZNF555\\
Mono / Ery & 10 & 0.119 & CDX2, IRF3, KLF6, MTERF3, NFYB, OVOL2, PRDM1, ZBTB49, ZNF300, ZNF555\\
\bottomrule\end{longtable}}
\clearpage\subsection{统计候选的文献证据与研究定位}
{\footnotesize\linespread{1}\selectfont\setlength{\tabcolsep}{4pt}\renewcommand{\arraystretch}{1.08}\begin{longtable}{p{2.6cm}rp{2.6cm}p{9cm}}\caption{统计证据—文献背景—验证方向。R 编号对应参考文献序号。}\\\toprule 目标 / TF & 排名 & 证据背景 & 文献依据与后续问题\\\midrule\endfirsthead\toprule 目标 / TF & 排名 & 证据背景 & 文献依据与后续问题\\\midrule\endhead
Erythrocyte / MYB & 1 & 红系机制支持 & 人 CD34+ 细胞中的 MYB 下调影响红系命运，并涉及 KLF1、LMO2；小鼠研究显示其作用随红系阶段变化。 将 MYB 相关匹配优先联系红系发育程序，检验阶段与表达持续时间。 R03、R04\\
Plasma cell / BCL6 & 15 & 浆细胞机制支持 & 小鼠 Bcl6 直接抑制 Prdm1；缺失后浆细胞发育加快。 负向匹配与解除 BCL6 抑制的方向相呼应；在 B 到浆细胞背景核对转变。 R05\\
Monocyte / KLF6 & 2 & 邻近髓系背景支持 & 人/鼠巨噬细胞研究支持 KLF6 的炎症调控；后续工作涉及 PRDM1–BCL6 轴。 优先考察炎症功能状态，并区分其与单核细胞谱系形成的关系。 R06、R07\\
Monocyte / PRDM1 & 8 & 邻近髓系背景支持 & 巨噬细胞 KLF6–PRDM1–BCL6 轴提供与 KLF6 共同出现的研究线索。 在对应背景检验两者共同响应与先后关系；本次未直接估计调控边。 R07\\
B cell / ZNF22 & 1 & 本次定向检索尚未定位直接证据 & 预定基因—B/淋巴/造血检索返回空；保留统计排名。 作为探索性候选，在独立 B 细胞数据中优先复核。 定向检索\\
Plasma cell / OVOL2 & 1 & 共同模式候选 & 本次定向检索命中一篇镉暴露/糖代谢研究，未作为血液分化机制证据。 比较其跨细胞共同表达程序与目标特有程序。 定向检索\\
Erythrocyte / NFYB & 3 & 共同模式候选 & 本阶段按四目标稳定交集讨论；未新增 NFYB 特异机制结论。 比较跨目标贡献，查明是否主要追踪共同状态特征。 定向检索\\
\bottomrule\end{longtable}}
文献核对以原始研究的摘要、可访问全文及官方索引为依据。部分全文访问受限，表中结论限定在已核对内容。BCL6 因机制证据列入讨论，未被追加到预设的前十自身读出删除实验；EBF1 作为已知 B 细胞参照保留其零系数结果。
{\footnotesize\linespread{1}\selectfont\setlength{\tabcolsep}{4pt}\renewcommand{\arraystretch}{1.08}\begin{longtable}{llr}\caption{相关扰动边（主阈值 $|r|\geq0.8$）。用于观察可替代模式。}\\\toprule 扰动 1 & 扰动 2 & Pearson\\\midrule\endfirsthead\toprule 扰动 1 & 扰动 2 & Pearson\\\midrule\endhead
PLA2G12A & SIRT3 & 0.809\\
HSP90AB1 & MYB & 0.801\\
POLB & DHH & 0.831\\
DHH & MYB & 0.802\\
DHH & RPS6KB2 & 0.824\\
PPP3CB & RPS6KB2 & 0.802\\
\bottomrule\end{longtable}}
\clearpage\subsection{Lasso 与 Elastic Net 的相关变量结构}
\plate{04_F5_correlated_predictor_groups}{c:correlation}{Lasso 与 Elastic Net 的全扰动系数对照及相关组。各点为扰动，重点前十 TF 高亮；连通组按相关阈值定义，保留组内不同成员。两种模型的活跃方向一致性需与候选组成和预测性能共同阅读。}{0.73}
{\footnotesize\linespread{1}\selectfont\setlength{\tabcolsep}{4pt}\renewcommand{\arraystretch}{1.08}\begin{longtable}{lrrrrr}\caption{全活跃扰动集合的一致性。}\\\toprule 目标 & Lasso 数 & EN 数 & 交集 & Jaccard & 系数 Spearman\\\midrule\endfirsthead\toprule 目标 & Lasso 数 & EN 数 & 交集 & Jaccard & 系数 Spearman\\\midrule\endhead
B & 656 & 656 & 655 & 0.997 & 0.999\\
Plasma & 784 & 836 & 760 & 0.884 & 0.956\\
Mono & 683 & 730 & 663 & 0.884 & 0.954\\
Ery & 593 & 700 & 566 & 0.779 & 0.902\\
\bottomrule\end{longtable}}
\clearpage\subsection{完整尺度与读出敏感性汇总}
原尺度差值、逐基因标准化差值和 landmark 子集具有各自响应与读出定义。表中保留各任务的原生评价尺度；同一 landmark 测试集上的训练覆盖比较在正文图 4D。等数量非 landmark 对照检验读出组成的影响。
{\footnotesize\linespread{1}\selectfont\setlength{\tabcolsep}{4pt}\renewcommand{\arraystretch}{1.08}\begin{longtable}{llrrrrr}\caption{全部尺度与训练读出比较。MSE 选择相同，列出两套 $R^2$ 与最低误差模型其他指标。}\\\toprule 任务 & 目标 & $R^2_{min}$ & $R^2_{1SE}$ & RMSE$_{min}$ & $r_{min}$ & $\rho_{min}$\\\midrule\endfirsthead\toprule 任务 & 目标 & $R^2_{min}$ & $R^2_{1SE}$ & RMSE$_{min}$ & $r_{min}$ & $\rho_{min}$\\\midrule\endhead
Raw & B & 0.427 & 0.397 & 0.1595 & 0.654 & 0.611\\
Raw & Plasma & 0.429 & 0.393 & 0.1620 & 0.655 & 0.583\\
Raw & Mono & 0.483 & 0.445 & 0.1927 & 0.696 & 0.628\\
Raw & Ery & 0.609 & 0.568 & 0.2413 & 0.781 & 0.695\\
Gene-scaled & B & 0.273 & 0.258 & 1.1080 & 0.522 & 0.562\\
Gene-scaled & Plasma & 0.307 & 0.280 & 1.1702 & 0.554 & 0.558\\
Gene-scaled & Mono & 0.312 & 0.291 & 1.1642 & 0.559 & 0.623\\
Gene-scaled & Ery & 0.413 & 0.394 & 1.0430 & 0.643 & 0.672\\
Landmark & B & 0.005 & -0.003 & 0.2669 & 0.077 & 0.113\\
Landmark & Plasma & 0.007 & -0.002 & 0.2661 & 0.086 & 0.108\\
Landmark & Mono & 0.034 & -0.001 & 0.3005 & 0.191 & 0.249\\
Landmark & Ery & -0.003 & -0.001 & 0.3495 & 0.043 & 0.108\\
All → landmark & B & 0.316 & 0.286 & 0.2213 & 0.592 & 0.576\\
All → landmark & Plasma & 0.324 & 0.274 & 0.2197 & 0.614 & 0.583\\
All → landmark & Mono & 0.328 & 0.302 & 0.2507 & 0.627 & 0.545\\
All → landmark & Ery & 0.224 & 0.200 & 0.3074 & 0.598 & 0.598\\
Size match 1 & B & 0.100 & -0.000 & 0.2053 & 0.317 & 0.352\\
Size match 1 & Plasma & 0.064 & -0.004 & 0.1963 & 0.264 & 0.258\\
Size match 1 & Mono & 0.226 & 0.070 & 0.2184 & 0.476 & 0.432\\
Size match 1 & Ery & 0.489 & 0.389 & 0.2884 & 0.700 & 0.557\\
Size match 2 & B & 0.093 & 0.007 & 0.2169 & 0.312 & 0.326\\
Size match 2 & Plasma & 0.179 & 0.057 & 0.1955 & 0.424 & 0.310\\
Size match 2 & Mono & 0.192 & 0.076 & 0.2162 & 0.439 & 0.387\\
Size match 2 & Ery & 0.361 & 0.241 & 0.3017 & 0.603 & 0.533\\
Size match 3 & B & 0.137 & 0.058 & 0.1784 & 0.372 & 0.350\\
Size match 3 & Plasma & 0.247 & 0.078 & 0.1842 & 0.498 & 0.397\\
Size match 3 & Mono & 0.358 & 0.211 & 0.2296 & 0.599 & 0.466\\
Size match 3 & Ery & 0.485 & 0.355 & 0.2909 & 0.697 & 0.533\\
\bottomrule\end{longtable}}
{\footnotesize\linespread{1}\selectfont\setlength{\tabcolsep}{4pt}\renewcommand{\arraystretch}{1.08}\begin{longtable}{lllll}\caption{前十候选的活跃方向保留与自身读出覆盖。}\\\toprule 目标 & Elastic Net & Gene-scaled & Landmark & 自身读出可匹配\\\midrule\endfirsthead\toprule 目标 & Elastic Net & Gene-scaled & Landmark & 自身读出可匹配\\\midrule\endhead
B & 10/10 & 9/10 & 0/10 & 7/10\\
Plasma & 10/10 & 9/10 & 0/10 & 7/10\\
Mono & 10/10 & 8/10 & 0/10 & 10/10\\
Ery & 10/10 & 6/10 & 0/10 & 10/10\\
\bottomrule\end{longtable}}
\clearpage\section{外部扰动实验}\label{app:D}\setcounter{figure}{0}\setcounter{table}{0}\renewcommand{\thefigure}{D\arabic{figure}}\renewcommand{\thetable}{D\arabic{table}}
本附录覆盖划分、训练、完整汇总与代表条件。统一输出 18,977 个共同基因，HVG1000 每次由训练集确定；单个条件等权，三个种子再等权。模型缩写 scF+GEARS 指冻结 scFoundation 编码器的任务模型。
{\small\linespread{1}\selectfont\setlength{\tabcolsep}{4pt}\renewcommand{\arraystretch}{1.08}\begin{longtable}{lrrrrrrr}\caption{条件划分与最终训练覆盖。各测试集另含 25 个未见单扰动。}\\\toprule 种子 & 训练 & 验证 & 测试 & Seen0 & Seen1 & Seen2 & 训练细胞\\\midrule\endfirsthead\toprule 种子 & 训练 & 验证 & 测试 & Seen0 & Seen1 & Seen2 & 训练细胞\\\midrule\endhead
1 & 98 & 18 & 109 & 18 & 56 & 10 & 44158\\
2 & 104 & 28 & 93 & 15 & 37 & 16 & 45661\\
3 & 104 & 24 & 97 & 10 & 48 & 14 & 45696\\
\bottomrule\end{longtable}}
{\footnotesize\linespread{1}\selectfont\setlength{\tabcolsep}{4pt}\renewcommand{\arraystretch}{1.08}\begin{longtable}{lrrrrrr}\caption{九项正式运行。轮数为实际记录；所有最佳检查点在前八轮。耗时含训练与预测，显存为最终进程的分配峰值。}\\\toprule 模型 & 种子 & 实际轮数 & 最佳轮 & 验证 MSE & 分钟 & 显存 GiB\\\midrule\endfirsthead\toprule 模型 & 种子 & 实际轮数 & 最佳轮 & 验证 MSE & 分钟 & 显存 GiB\\\midrule\endhead
GEARS & 1 & 12 & 7 & 0.00162 & 20.1 & 6.1\\
GEARS & 2 & 10 & 5 & 0.00184 & 17.4 & 6.1\\
GEARS & 3 & 8 & 5 & 0.00197 & 13.9 & 6.1\\
scGPT & 1 & 8 & 8 & 0.00132 & 90.9 & 5.3\\
scGPT & 2 & 8 & 7 & 0.00137 & 94.2 & 5.3\\
scGPT & 3 & 8 & 6 & 0.00168 & 89.4 & 15.9\\
scF+GEARS & 1 & 9 & 8 & 0.00202 & 72.5 & 7.8\\
scF+GEARS & 2 & 8 & 4 & 0.00233 & 60.9 & 7.8\\
scF+GEARS & 3 & 6 & 1 & 0.00276 & 45.5 & 20.7\\
\bottomrule\end{longtable}}
\plate{07_F3_training_curves}{d:training}{所有模型与种子的验证曲线。点标最佳检查点，灰区标示预算变更前已完成的第八轮之后历史；这些额外轮次没有进入最终检查点选择。}{0.43}
\clearpage\subsection{总体指标与直接靶基因排除}
{\footnotesize\linespread{1}\selectfont\setlength{\tabcolsep}{4pt}\renewcommand{\arraystretch}{1.08}\begin{longtable}{lllrrrrr}\caption{全部方法的总体指标。SD 为三个划分间的标准差；相关无定义记为 ---。}\\\toprule 直接靶基因 & 读出 & 模型 & MSE & SD & RMSE & $r$ & $\rho$\\\midrule\endfirsthead\toprule 直接靶基因 & 读出 & 模型 & MSE & SD & RMSE & $r$ & $\rho$\\\midrule\endhead
包含 & Shared & Zero & 0.00304 & 0.00026 & 0.0551 & --- & ---\\
包含 & Shared & Train mean & 0.00197 & 0.00026 & 0.0443 & 0.596 & 0.398\\
包含 & Shared & Ridge add. & 0.00197 & 0.00035 & 0.0443 & 0.701 & 0.423\\
包含 & Shared & GEARS & 0.00221 & 0.00030 & 0.0470 & 0.467 & 0.221\\
包含 & Shared & scGPT & 0.00203 & 0.00030 & 0.0450 & 0.544 & 0.311\\
包含 & Shared & scF+GEARS & 0.00275 & 0.00014 & 0.0525 & 0.422 & 0.143\\
包含 & HVG & Zero & 0.03412 & 0.00247 & 0.1846 & --- & ---\\
包含 & HVG & Train mean & 0.02051 & 0.00248 & 0.1430 & 0.587 & 0.610\\
包含 & HVG & Ridge add. & 0.02091 & 0.00353 & 0.1443 & 0.739 & 0.694\\
包含 & HVG & GEARS & 0.01974 & 0.00245 & 0.1403 & 0.547 & 0.514\\
包含 & HVG & scGPT & 0.02068 & 0.00304 & 0.1435 & 0.546 & 0.524\\
包含 & HVG & scF+GEARS & 0.02340 & 0.00123 & 0.1529 & 0.490 & 0.477\\
排除 & Shared & Zero & 0.00300 & 0.00026 & 0.0548 & --- & ---\\
排除 & Shared & Train mean & 0.00193 & 0.00026 & 0.0439 & 0.601 & 0.398\\
排除 & Shared & Ridge add. & 0.00195 & 0.00035 & 0.0440 & 0.698 & 0.423\\
排除 & Shared & GEARS & 0.00218 & 0.00030 & 0.0466 & 0.470 & 0.221\\
排除 & Shared & scGPT & 0.00202 & 0.00030 & 0.0449 & 0.534 & 0.311\\
排除 & Shared & scF+GEARS & 0.00272 & 0.00015 & 0.0521 & 0.424 & 0.143\\
排除 & HVG & Zero & 0.03404 & 0.00247 & 0.1844 & --- & ---\\
排除 & HVG & Train mean & 0.02042 & 0.00248 & 0.1427 & 0.588 & 0.610\\
排除 & HVG & Ridge add. & 0.02089 & 0.00352 & 0.1442 & 0.736 & 0.694\\
排除 & HVG & GEARS & 0.01965 & 0.00244 & 0.1400 & 0.548 & 0.514\\
排除 & HVG & scGPT & 0.02067 & 0.00303 & 0.1435 & 0.541 & 0.524\\
排除 & HVG & scF+GEARS & 0.02331 & 0.00124 & 0.1526 & 0.492 & 0.477\\
\bottomrule\end{longtable}}
三次划分不是独立生物重复，SD 描述划分差异。相关系数按有定义的条件汇总。对零变化或全未见组成下的恒定预测，相关系数无定义，但 MSE／RMSE 仍可计算。
\clearpage\subsection{泛化分层的完整误差与模式相关性}
{\footnotesize\linespread{1}\selectfont\setlength{\tabcolsep}{4pt}\renewcommand{\arraystretch}{1.08}\begin{longtable}{lllrrrr}\caption{分层评价保留全部方法。Seen0/1/2 按最终训练条件重新核对。}\\\toprule 读出 & 条件组 & 模型 & MSE & RMSE & $r$ & $\rho$\\\midrule\endfirsthead\toprule 读出 & 条件组 & 模型 & MSE & RMSE & $r$ & $\rho$\\\midrule\endhead
Shared & Single & Zero & 0.00186 & 0.0430 & --- & ---\\
Shared & Single & Train mean & 0.00126 & 0.0354 & 0.559 & 0.371\\
Shared & Single & Ridge add. & 0.00186 & 0.0430 & --- & ---\\
Shared & Single & GEARS & 0.00164 & 0.0404 & 0.334 & 0.168\\
Shared & Single & scGPT & 0.00151 & 0.0387 & 0.417 & 0.244\\
Shared & Single & scF+GEARS & 0.00199 & 0.0445 & 0.365 & 0.122\\
Shared & seen0 & Zero & 0.00377 & 0.0613 & --- & ---\\
Shared & seen0 & Train mean & 0.00229 & 0.0475 & 0.673 & 0.446\\
Shared & seen0 & Ridge add. & 0.00377 & 0.0613 & --- & ---\\
Shared & seen0 & GEARS & 0.00281 & 0.0528 & 0.534 & 0.249\\
Shared & seen0 & scGPT & 0.00270 & 0.0517 & 0.566 & 0.331\\
Shared & seen0 & scF+GEARS & 0.00319 & 0.0564 & 0.485 & 0.158\\
Shared & seen1 & Zero & 0.00336 & 0.0579 & --- & ---\\
Shared & seen1 & Train mean & 0.00218 & 0.0466 & 0.598 & 0.401\\
Shared & seen1 & Ridge add. & 0.00176 & 0.0419 & 0.658 & 0.396\\
Shared & seen1 & GEARS & 0.00232 & 0.0481 & 0.504 & 0.233\\
Shared & seen1 & scGPT & 0.00214 & 0.0461 & 0.578 & 0.330\\
Shared & seen1 & scF+GEARS & 0.00299 & 0.0547 & 0.431 & 0.146\\
Shared & seen2 & Zero & 0.00317 & 0.0562 & --- & ---\\
Shared & seen2 & Train mean & 0.00202 & 0.0447 & 0.598 & 0.409\\
Shared & seen2 & Ridge add. & 0.00084 & 0.0288 & 0.848 & 0.511\\
Shared & seen2 & GEARS & 0.00205 & 0.0448 & 0.525 & 0.257\\
Shared & seen2 & scGPT & 0.00173 & 0.0415 & 0.659 & 0.358\\
Shared & seen2 & scF+GEARS & 0.00270 & 0.0518 & 0.434 & 0.166\\
HVG & Single & Zero & 0.02072 & 0.1435 & --- & ---\\
HVG & Single & Train mean & 0.01304 & 0.1138 & 0.533 & 0.562\\
HVG & Single & Ridge add. & 0.02072 & 0.1435 & --- & ---\\
HVG & Single & GEARS & 0.01539 & 0.1238 & 0.399 & 0.370\\
HVG & Single & scGPT & 0.01551 & 0.1241 & 0.380 & 0.380\\
HVG & Single & scF+GEARS & 0.01611 & 0.1267 & 0.429 & 0.429\\
HVG & seen0 & Zero & 0.04301 & 0.2070 & --- & ---\\
HVG & seen0 & Train mean & 0.02405 & 0.1543 & 0.675 & 0.690\\
HVG & seen0 & Ridge add. & 0.04301 & 0.2070 & --- & ---\\
HVG & seen0 & GEARS & 0.02611 & 0.1607 & 0.627 & 0.597\\
HVG & seen0 & scGPT & 0.02872 & 0.1689 & 0.567 & 0.565\\
HVG & seen0 & scF+GEARS & 0.02710 & 0.1645 & 0.565 & 0.540\\
HVG & seen1 & Zero & 0.03763 & 0.1939 & --- & ---\\
HVG & seen1 & Train mean & 0.02266 & 0.1503 & 0.589 & 0.612\\
HVG & seen1 & Ridge add. & 0.01789 & 0.1334 & 0.696 & 0.653\\
HVG & seen1 & GEARS & 0.02013 & 0.1418 & 0.584 & 0.550\\
HVG & seen1 & scGPT & 0.02155 & 0.1463 & 0.590 & 0.561\\
HVG & seen1 & scF+GEARS & 0.02585 & 0.1606 & 0.497 & 0.481\\
HVG & seen2 & Zero & 0.03524 & 0.1874 & --- & ---\\
HVG & seen2 & Train mean & 0.02091 & 0.1440 & 0.602 & 0.633\\
HVG & seen2 & Ridge add. & 0.00687 & 0.0827 & 0.892 & 0.829\\
HVG & seen2 & GEARS & 0.01735 & 0.1297 & 0.615 & 0.581\\
HVG & seen2 & scGPT & 0.01647 & 0.1281 & 0.704 & 0.635\\
HVG & seen2 & scF+GEARS & 0.02268 & 0.1499 & 0.508 & 0.498\\
\bottomrule\end{longtable}}
\clearpage\subsection{直接加和的匹配比较与权重来源}
\plate{07_F4_direct_additive_subset}{d:direct}{直接单扰动加和的可用子集。三个种子分别有 10、16、14 个可匹配组合；全部七种方法在相同条件子集内评价。比值相对直接加和计算，点与误差线保留划分间波动。}{0.43}
{\footnotesize\linespread{1}\selectfont\setlength{\tabcolsep}{4pt}\renewcommand{\arraystretch}{1.08}\begin{longtable}{llrrrr}\caption{匹配子集完整指标。}\\\toprule 模型 & 读出 & MSE & RMSE & $r$ & $\rho$\\\midrule\endfirsthead\toprule 模型 & 读出 & MSE & RMSE & $r$ & $\rho$\\\midrule\endhead
Zero & Shared & 0.00317 & 0.0562 & --- & ---\\
Zero & HVG & 0.03524 & 0.1874 & --- & ---\\
Train mean & Shared & 0.00202 & 0.0447 & 0.598 & 0.409\\
Train mean & HVG & 0.02091 & 0.1440 & 0.602 & 0.633\\
Ridge add. & Shared & 0.00084 & 0.0288 & 0.848 & 0.511\\
Ridge add. & HVG & 0.00687 & 0.0827 & 0.892 & 0.829\\
GEARS & Shared & 0.00205 & 0.0448 & 0.525 & 0.257\\
GEARS & HVG & 0.01735 & 0.1297 & 0.615 & 0.581\\
scGPT & Shared & 0.00173 & 0.0415 & 0.659 & 0.358\\
scGPT & HVG & 0.01647 & 0.1281 & 0.704 & 0.635\\
scF+GEARS & Shared & 0.00270 & 0.0518 & 0.434 & 0.166\\
scF+GEARS & HVG & 0.02268 & 0.1499 & 0.508 & 0.498\\
Direct add. & Shared & 0.00109 & 0.0330 & 0.839 & 0.513\\
Direct add. & HVG & 0.00955 & 0.0976 & 0.884 & 0.824\\
\bottomrule\end{longtable}}
{\footnotesize 源码版本：GEARS \texttt{f374e43e197b}；scGPT \texttt{cebd6fae655b}；scFoundation \texttt{397631c495ed}。scGPT 使用官方 blood 通用权重，scFoundation 使用公共镜像固定版本 \texttt{cb434153a1ac}，文件 SHA-256 与该镜像发布标识一致。运行环境为 Python 3.12.13、PyTorch 2.12.0+cu130。预训练数据与 Norman 的样本重叠状态未核实。}\par
\clearpage\subsection{不同响应幅度的完整代表预测}
\plate{07_F5_representative_predictions}{d:cases}{种子 1 测试组合中，按真实 HVG 响应均方根选取最大、中位和最小响应条件。每行同一条件，保留 Ridge、GEARS、scGPT 与 scFoundation＋GEARS 全部预测；同一行采用共同坐标范围。该选择覆盖不同信号强度，不依据模型误差筛选。}{0.84}
\clearpage\subsection{PRDM1 的完整响应及指标}
\plate{Supp_PRDM1}{d:prdm}{PRDM1 衔接案例的完整方法预测。上排单扰动，下排 CBFA2T3＋PRDM1 组合；每种方法均显示同一套训练 HVG1000。各排共享坐标范围，虚线为一致线。图中使用连续 A–L 标识。}{0.42}
{\footnotesize\linespread{1}\selectfont\setlength{\tabcolsep}{4pt}\renewcommand{\arraystretch}{1.08}\begin{longtable}{llrrrr}\caption{PRDM1 及其组合的所有方法。种子 1；改进为相对零变化 MSE 降低百分比。}\\\toprule 条件 & 模型 & MSE & $r$ & $\rho$ & 改进\%\\\midrule\endfirsthead\toprule 条件 & 模型 & MSE & $r$ & $\rho$ & 改进\%\\\midrule\endhead
单扰动 & Zero & 0.04378 & --- & --- & 0.0\\
单扰动 & Train mean & 0.02107 & 0.801 & 0.784 & 51.9\\
单扰动 & Ridge add. & 0.04378 & --- & --- & 0.0\\
单扰动 & GEARS & 0.02046 & 0.821 & 0.699 & 53.3\\
单扰动 & scGPT & 0.03121 & 0.615 & 0.594 & 28.7\\
单扰动 & scF+GEARS & 0.01710 & 0.751 & 0.633 & 60.9\\
组合 & Zero & 0.08208 & --- & --- & 0.0\\
组合 & Train mean & 0.04351 & 0.885 & 0.878 & 47.0\\
组合 & Ridge add. & 0.08208 & --- & --- & 0.0\\
组合 & GEARS & 0.03009 & 0.861 & 0.845 & 63.3\\
组合 & scGPT & 0.04725 & 0.809 & 0.800 & 42.4\\
组合 & scF+GEARS & 0.02901 & 0.768 & 0.724 & 64.7\\
\bottomrule\end{longtable}}
核心稳定 TF 与 Norman 靶标交集为 EGR1、OSR2、PRDM1、HNF4A、ZBTB25，共八个目标细胞—候选条目。PRDM1 单扰动在种子 1 为未见单扰动，组合为 seen0；在种子 2、3 中，PRDM1 单扰动处于训练集。直接单扰动加和不覆盖这两个种子 1 条件。

\end{document}
